\documentclass{beamer}
\usepackage[T1]{fontenc}
\usepackage[utf8]{inputenc}
\usepackage{lmodern}
\usepackage{amsmath,amssymb,mathtools,bm}
\usepackage{booktabs}
\usetheme{Boadilla}
\useoutertheme{miniframes}
\setbeamertemplate{navigation symbols}{}
\newcommand{\dd}{\mathrm{d}}
\newcommand{\ii}{\mathrm{i}}
\newcommand{\norm}[1]{\left\|#1\right\|}


\title{MATH4907 Tutorial 3}
\author{Sanyang Liu}
\date{October 25, 2026}

\begin{document}
\maketitle
\AtBeginSection[]
{
    \begin{frame}{Table of Contents}
        \tableofcontents[
            currentsection,
            hideallsubsections
        ]
    \end{frame}
}

\begin{frame}{Personal Information}
\begin{itemize}
  \item TA: Sanyang Liu
  \item Office: Room 320A Run Run Shaw Building (Email me first)
  \item Email: \texttt{u3013241@connect.hku.hk}
\end{itemize}
\end{frame}

\begin{frame}{Convergence of Scheme}
Last tutorial, we studied
\begin{itemize}
    \item Finite difference stencils and their order of accuracy
    \item Finite difference schemes and their local truncation errors
\end{itemize}

However, local accuracy alone \textbf{does not guarantee} global convergence of the numerical solution to the exact solution as the grid is refined.

\begin{block}{Definition: Convergence of a Scheme}
Let $u^j$ denote the exact solution sampled on the grid at time $\tau_j$,
and let $v^j$ be the numerical solution. The scheme is (conditionally)
convergent with respect to the norm $\|\cdot\|$ if, for every fixed $t>0$,
\[
\|v^j-u^j\|\longrightarrow 0
\qquad\text{as } h,k\longrightarrow 0^+,
\]
whenever $(h,k)\in\Lambda$ and $\tau_j\longrightarrow t$, for some
admissible set $\Lambda\subset(0,\infty)^2$.
\end{block}
\end{frame}


\section{Consistency}
\begin{frame}{Order of Accuracy of the $\theta$-Scheme}
\begin{block}{Question 1}
Consider the heat equation $u_\tau-u_{xx}=0$
and the $\theta$-scheme
\[
v_{j+1}=v_j-(1-\theta)\lambda Bv_j-\theta\lambda Bv_{j+1},
\qquad \theta\in(0,1).
\]
Determine the order of accuracy of this scheme.
\end{block}
\end{frame}

\begin{frame}{Substituting Exact Solution}
\begin{block}{$\theta$-Scheme}
\[
\frac{v_{i,j+1}-v_{i,j}}{k}
=
(1-\theta)\frac{v_{i-1,j}-2v_{i,j}+v_{i+1,j}}{h^2}
+\theta\frac{v_{i-1,j+1}-2v_{i,j+1}+v_{i+1,j+1}}{h^2},
\]
where $x_i=ih$ and $\tau_j=jk$.
\end{block}
Substituting the exact solution into the scheme and using Taylor
expansions gives the following approximations:
\[
\frac{u(x_i,\tau_{j+1})-u(x_i,\tau_j)}{k}
=
u_\tau(x_i,\tau_j)
+\frac{k}{2}u_{\tau\tau}(x_i,\tau_j)
+O(k^2).
\]
\[
\frac{u(x_{i-1},\tau_j)-2u(x_i,\tau_j)+u(x_{i+1},\tau_j)}{h^2}
=
u_{xx}(x_i,\tau_j)+O(h^2).
\]
\end{frame}

\begin{frame}{Local Truncation Error}
At $\tau_{j+1}$, the centered difference gives
\[
\frac{u(x_{i-1},\tau_{j+1})-2u(x_i,\tau_{j+1})+u(x_{i+1},\tau_{j+1})}{h^2}
=
u_{xx}(x_i,\tau_{j+1})+O(h^2).
\]
Expanding $u_{xx}$ in time about $\tau_j$,
\[
u_{xx}(x_i,\tau_{j+1})
=
u_{xx}(x_i,\tau_j)+k u_{xx\tau}(x_i,\tau_j)+O(k^2).
\]

Therefore, the local truncation error is
\[
\begin{aligned}
\mathcal{T}_{i,j}
&=
u_\tau-(1-\theta)u_{xx}-\theta u_{xx}
+k\left(\frac12u_{\tau\tau}-\theta u_{xx\tau}\right)
+O(h^2)+O(k^2)\\
&=
{u_\tau-u_{xx}}
+k\left(\frac12u_{\tau\tau}-\theta u_{xx\tau}\right)
+O(h^2)+O(k^2).
\end{aligned}
\]
\end{frame}

\begin{frame}{Order of Accuracy}
Using the heat equation and differentiating it with respect to $\tau$,
\[
u_\tau=u_{xx},
\qquad
u_{\tau\tau}=u_{xx\tau},
\]
the local truncation error becomes
\[
\mathcal{T}_{i,j}
=
\left(\frac12-\theta\right)k u_{\tau\tau}
+O(h^2)+O(k^2).
\]

If $\theta\neq\frac12$, then
\[
\boxed{\mathcal{T}_{i,j}=O(h^2)+O(k),}
\]
If $\theta=\frac12$, the first-order time term vanishes, giving
\[
\boxed{\mathcal{T}_{i,j}=O(h^2)+O(k^2).}
\]
Thus, the Crank--Nicolson scheme is second-order accurate in both space and time.

\end{frame}

\begin{frame}{Consistency}

\begin{block}{Definition: Consistency of Scheme}
Assume the scheme can be written as $v^{j+1}=Q(v^j)+k\,G^j,$
where $Q:V\to V$ is linear and $G^j$ encodes the boundary conditions.

Let $u^j$ be the exact solution sampled on the grid. The scheme is
\textbf{consistent} with the PDE in the norm $\|\cdot\|$ if
\[
\left\|
\frac{u^{j+1}-Q(u^j)}{k}-G^j
\right\|\longrightarrow 0
\quad\text{as }h,k\longrightarrow0^+,
\qquad (h,k)\in\Lambda,
\]
for an admissible set $\Lambda\subset(0,\infty)^2$.
\end{block}
\begin{block}{Remark}
The norm $\|\cdot\|$ may depend on $h=\Delta x$,
since the length of the grid vectors changes with the grid (e.g., $\|x\|_{2,\Delta x_n}^2
=
\Delta x_n\|x\|_2^2
=
\Delta x_n\sum_{i=1}^n |x_i|^2$).
\end{block}
\end{frame}

\begin{frame}{Consistency of the $\theta$-Scheme}

For homogeneous boundary conditions, write the scheme as
\[
v^{j+1}=Q(v^j),\qquad
Q=(I+\theta\lambda B)^{-1}
\bigl(I-(1-\theta)\lambda B\bigr),
\qquad \lambda=\frac{k}{h^2}.
\]

The truncation error is
\[
\boldsymbol{\tau}^j=\frac{1}{k}\bigl(u^{j+1}-Q(u^j)\bigr),
\]
whose $i$th component is the local truncation error
\[
\mathcal{T}_{i,j}
=
\left(\frac12-\theta\right)k\,u_{\tau\tau}(x_i,\tau_j)
+O(h^2)+O(k^2).
\]

Thus, $\|\boldsymbol{\tau}^j\|\to0$ as $h,k\to0^+$, so the scheme is consistent.

\end{frame}

\begin{frame}{Consistency Alone Does Not Ensure Convergence}

For $u_\tau=u_{xx}$ with $u(x,0)=\sin(\pi x)$,
\[
u(x,\tau)=e^{-\pi^2\tau}\sin(\pi x),
\qquad
u(0.2,0.5)\approx0.004227.
\]

The explicit scheme has local truncation error
$O(k+h^2)$, but its stability depends on $\lambda=k/h^2$.

\begin{center}
\begin{tabular}{cccc}
\toprule
$h$ & $k$ & $\lambda=k/h^2$ & Numerical value at $(0.2,0.5)$\\
\midrule
$0.1$ & $0.0005$ & $0.05$ & $0.0043492$\\
$0.1$ & $0.01$ & $1$ & $-2694895.2$\\
\bottomrule
\end{tabular}
\end{center}

When $\lambda=1$, errors grow as the scheme advances in time.
To describe how errors propagate and whether they remain bounded,
we introduce the concept of \textbf{stability} of the scheme.
\end{frame}



\section{Stability}
\begin{frame}{Stability}
We still consider a two-level finite difference scheme
\[
v^{j+1}=Q(v^j).
\]
The scheme is \textbf{stable} with respect to the norm $\|\cdot\|$ if there
exist constants $\Delta t_0,\Delta x_0,K>0$ and $\beta\ge 0$ such that
\[
\|v^j\|\le K e^{\beta j\Delta t}\|v^0\|
\]
whenever
\[
(\Delta x,\Delta t)\in
\bigl((0,\Delta x_0)\times(0,\Delta t_0)\bigr)\cap\Lambda,
\qquad \Lambda\subset(0,\infty)^2.
\]
\begin{block}{Remark}
Stability of a finite difference scheme is analogous to well-posedness
of the PDE.
\end{block}

\end{frame}

\begin{frame}{Lax Equivalence Theorem}
\begin{block}{Theorem: Lax Equivalence Theorem}
For a well-posed initial-value problem for a linear PDE, a consistent
two-level finite difference scheme is convergent \textbf{if and only if} it is
stable.
\end{block}

\begin{itemize}
    \item A well-posed IVP has a solution that depends continuously on its
    initial data.
    \item A two-level scheme involves only two consecutive time levels,
    $t_j$ and $t_{j+1}$.
\end{itemize}

\vfill
{\footnotesize
Reference: P.~D.~Lax and R.~D.~Richtmyer,
``Survey of the Stability of Linear Finite Difference Equations,''
\textit{Communications on Pure and Applied Mathematics},
vol.~9, no.~2, pp.~267--293, 1956.
}

\end{frame}






\begin{frame}{Discrete Fourier Transform}

\begin{block}{Definition: Discrete Fourier Transform}
\[
\mathcal{F}:\ell^2(\mathbb{Z})\longrightarrow L^2([-\pi,\pi]),
\qquad
(\mathcal{F}v)(\xi)=\widehat v(\xi)
=\frac{1}{\sqrt{2\pi}}\sum_{j\in\mathbb Z}v_j e^{-\ii j\xi}.
\]
\end{block}
Here, $\ell^2(\mathbb{Z})$ is the space of square-summable sequences,
and $L^2([-\pi,\pi])$ is the space of square-integrable functions.
\begin{block}{Theorem: Parseval's Identity}
For every $v\in\ell^2(\mathbb{Z})$,
\[
\|\mathcal{F}v\|_{L^2([-\pi,\pi])}^2
=
\int_{-\pi}^{\pi}|\widehat v(\xi)|^2\,\dd\xi
=
\sum_{j\in\mathbb Z}|v_j|^2
=
\|v\|_{\ell^2(\mathbb{Z})}^2.
\]
Thus, the Fourier transform preserves the norm.
\end{block}
\end{frame}


\begin{frame}{Fourier Transform of a Basis Vector}

\begin{block}{Question 2}
Let $e_j=(0,\ldots,1,\ldots,0,\ldots)$ be the $j$th canonical basis
vector in $\ell^2$, such that only the $j$th entry is nonzero and equals
to $1$. Then compute $\widehat{e_j}$.
\end{block}

The entries of $e_j$ are
\[
(e_j)_k=\delta_{jk}
=
\begin{cases}
1, & k=j,\\
0, & k\ne j.
\end{cases}
\]
By the definition of the discrete Fourier transform,
\[
\widehat{e_j}(\xi)
=
\frac{1}{\sqrt{2\pi}}
\sum_{k\in\mathbb{Z}}(e_j)_k e^{-\ii k\xi}
=
\frac{e^{-\ii j\xi}}{\sqrt{2\pi}}.
\]

\end{frame}



\begin{frame}{Why Use Fourier Analysis?}

Let $v^j=(v_i^j)_{i\in\mathbb{Z}}$ denote the numerical solution at time
level $j$. For $p\in\mathbb{Z}$, define the shift operator
\[
(\tau_p v^j)_i=v_{i+p}^j.
\]
A linear, constant-coefficient scheme can be written as
\[
v^{j+1}=\sum_p c_p\,\tau_p(v^j).
\]
The Fourier transform converts each shift into multiplication:
\[
\mathcal{F}(\tau_p v^j)(\xi)
=e^{\ii p\xi}\widehat v^{\,j}(\xi),
\qquad
\widehat v^{\,j}=\mathcal{F}(v^j).
\]
Therefore,
\[
\widehat v^{\,j+1}(\xi)
=g(\xi)\widehat v^{\,j}(\xi),
\qquad
g(\xi)=\sum_p c_p e^{\ii p\xi}.
\]
Each frequency evolves independently. By Parseval's identity, controlling $|g(\xi)|^j$ for all frequencies
controls the $\ell^2$ norm of the solution.

\end{frame}

\begin{frame}{Fourier Symbols of Difference Stencils}

For the Fourier mode $v_i=e^{\ii i\xi}$, the first-difference symbols are
\[
\begin{aligned}
\delta_x^+ = \frac{v_{i+1}-v_i}{h}
&\longleftrightarrow
\frac{e^{\ii\xi}-1}{h}
=\frac{\cos\xi-1+\ii\sin\xi}{h}
=\frac{-2\sin^2(\xi/2)+\ii\sin\xi}{h},\\[2pt]
\delta_x^- = \frac{v_{i}-v_{i-1}}{h}
&\longleftrightarrow
\frac{1-e^{-\ii\xi}}{h}
=\frac{1-\cos\xi+\ii\sin\xi}{h}
=\frac{2\sin^2(\xi/2)+\ii\sin\xi}{h},\\[2pt]
\delta_x^0 = \frac{v_{i+1}-v_{i-1}}{2h}
&\longleftrightarrow
\frac{e^{\ii\xi}-e^{-\ii\xi}}{2h}
=\frac{\ii\sin\xi}{h}.
\end{aligned}
\]
The centered second-difference symbol is
\[
\delta_{xx}=\frac{v_{i-1}-2v_i+v_{i+1}}{h^2}
\longleftrightarrow
\frac{e^{-\ii\xi}-2+e^{\ii\xi}}{h^2}
=
\frac{2\cos\xi-2}{h^2}
=
-\frac{4\sin^2(\xi/2)}{h^2}.
\]

\end{frame}

\begin{frame}{Stability and the Amplification Factor}

Recall that the scheme is stable in $\|\cdot\|_{2,\Delta x}$ if there exist
$K,\Delta x_0,\Delta t_0>0$ and $\beta\ge0$ such that
\[
\|v^j\|_{2,\Delta x}
\le K e^{\beta j\Delta t}\|v^0\|_{2,\Delta x}
\]
for all admissible mesh sizes.

\begin{block}{Theorem: Stability and Amplification Factor}
Let $g(\xi)$ be the amplification factor of the scheme. The scheme is
stable in $\|\cdot\|_{2,\Delta x}$ if and only if there exist
$K,\Delta x_0,\Delta t_0>0$ and $\beta\ge0$ such that
\[
|g(\xi)|^j\le K e^{\beta j\Delta t},
\qquad \forall\,\xi\in[-\pi,\pi],
\]
whenever $0<\Delta x\le\Delta x_0$ and
$0<\Delta t\le\Delta t_0$.
\end{block}
\end{frame}


\begin{frame}{Amplification Bound Implies Stability}
Let $g(\xi)$ be the amplification factor. Fourier transformation gives
\[
\widehat v^{\,j+1}(\xi)
=g(\xi)\widehat v^{\,j}(\xi),
\qquad
\widehat v^{\,j}(\xi)
=g(\xi)^j\widehat v^{\,0}(\xi).
\]
Suppose that, for all $\xi\in[-\pi,\pi]$,
\[
|g(\xi)|^j\le K e^{\beta jk}.
\]
By Parseval's identity,
\[
\begin{aligned}
\|v^j\|_2^2
=\|\widehat v^{\,j}\|_2^2\le K^2 e^{2\beta jk}\|\widehat v^{\,0}\|_2^2=K^2 e^{2\beta jk}\|v^0\|_2^2.
\end{aligned}
\]
Therefore,
\[
\|v^j\|_2\le K e^{\beta jk}\|v^0\|_2,
\]
so the scheme is stable.

\end{frame}


\begin{frame}{Stability Implies the Amplification Bound}

Assume the scheme is stable:
\[
\|v^j\|_2\le K e^{\beta jk}\|v^0\|_2.
\]
Using the Fourier representation and Parseval's identity,
\[
\int_{-\pi}^{\pi}|g(\xi)|^{2j}|\widehat v^{\,0}(\xi)|^2\,\dd\xi
\le
K^2e^{2\beta jk}
\int_{-\pi}^{\pi}|\widehat v^{\,0}(\xi)|^2\,\dd\xi.
\]

Since $\widehat v^{\,0}$ can be any function in $L^2([-\pi,\pi])$,
\[
|g(\xi)|^j\le K e^{\beta jk}
\quad\text{for almost every }\xi\in[-\pi,\pi].
\]
Because $g$ is continuous, the bound holds for every $\xi$:
if it failed at some point, continuity would make it fail on an
interval of positive length, contradicting the almost-everywhere bound.
\end{frame}








\section[Parabolic]{Parabolic Equation}
\begin{frame}{Question 4}
\begin{block}{Parabolic Problem}
Consider
\[
u_t=u_{xx}-u_x,\qquad 0<x<1,\quad 0<t\le1,
\]
with boundary and initial conditions
\[
u(0,t)=u(1,t)=0,\qquad 0\le t\le1,
\qquad
u(x,0)=f(x),\quad 0<x<1.
\]
\end{block}
Let $x_i=ih$, $t_j=j\Delta t$, with $h=1/n$ and $\Delta t=1/m$.
The scheme uses Crank--Nicolson for diffusion and an explicit centered
difference for the advection term:
\[
\begin{aligned}
U_i^{j+1}-U_i^j
={}&\frac{\Delta t}{2h^2}
\Big[
U_{i-1}^{j+1}-2U_i^{j+1}+U_{i+1}^{j+1}
+U_{i-1}^{j}-2U_i^{j}+U_{i+1}^{j}
\Big]\\
&-\frac{\Delta t}{2h}
\Big[
U_{i+1}^{j}-U_{i-1}^{j}
\Big],
\end{aligned}
\]
for $j=0,\ldots,m-1$ and $i=1,\ldots,n-1$.

\end{frame}


\begin{frame}{Q4(a): Matrix Form}
Let $U^j=(U_1^j,\ldots,U_{n-1}^j)^\mathsf{T}$. Define
\[
B=\frac{1}{h^2}
\begin{pmatrix}
-2&1\\
1&-2&1\\
&\ddots&\ddots&\ddots\\
&&1&-2&1\\
&&&1&-2
\end{pmatrix},
\qquad
C=\frac{1}{2h}
\begin{pmatrix}
0&1\\
-1&0&1\\
&\ddots&\ddots&\ddots\\
&&-1&0&1\\
&&&-1&0
\end{pmatrix}.
\]
These represent the centered difference approximations of $u_{xx}$ and
$u_x$, respectively. Since the boundary values are zero, the scheme is
\[
\frac{U^{j+1}-U^j}{k}
=
\frac12B(U^{j+1}+U^j)-CU^j.
\]

Collecting the new-time terms gives
\[
\boxed{
\left(I-\frac{k}{2}B\right)U^{j+1}
=
\left(I+\frac{k}{2}B-kC\right)U^j.
}
\]
\end{frame}

\begin{frame}{Q4(b): Local Truncation Error}
Assume
\[
u\in C^{4,3}\bigl([0,1]\times[0,1]\bigr),
\qquad
u_{i,j}:=u(x_i,t_j).
\]
Define the centered spatial differences
\[
\delta_x u_{i,j}
:=\frac{u_{i+1,j}-u_{i-1,j}}{2h},
\qquad
\delta_{xx}u_{i,j}
:=\frac{u_{i-1,j}-2u_{i,j}+u_{i+1,j}}{h^2}.
\]

The local truncation error at $(x_i,t_j)$ is
\[
\mathcal{T}_{i,j}
=
\frac{u_{i,j+1}-u_{i,j}}{\Delta t}
-\frac12\left(\delta_{xx}u_{i,j+1}+\delta_{xx}u_{i,j}\right)
+\delta_xu_{i,j}.
\]
    
\end{frame}


\begin{frame}{Q4(b): Local Truncation Error (cont'd)}
Taylor expansion gives
\[
\frac{u_{i,j+1}-u_{i,j}}{k}
=
u_\tau+\frac{k}{2}u_{\tau\tau}+O(k^2),
\]
\[
\frac12\left(\delta_{xx}u_{i,j+1}+\delta_{xx}u_{i,j}\right)
=
u_{xx}+\frac{k}{2}u_{xx\tau}+O(h^2)+O(k^2),
\]
and
\[
\delta_xu_{i,j}=u_x+O(h^2),
\]
where the derivatives are evaluated at $(x_i,\tau_j)$. Hence,
\[
\mathcal{T}_{i,j}
=
u_\tau-u_{xx}+u_x
+\frac{k}{2}\left(u_{\tau\tau}-u_{xx\tau}\right)
+O(h^2)+O(k^2).
\]
\end{frame}

\begin{frame}{Q4(b): Local Truncation Error (cont'd)}
Using $u_\tau=u_{xx}-u_x$ and differentiating in $\tau$,
\[
u_{\tau\tau}=u_{xx\tau}-u_{x\tau}.
\]
Therefore,
\[
\mathcal{T}_{i,j}
=
-\frac{k}{2}u_{x\tau}+O(h^2)+O(k^2)
=
O(h^2)+O(k).
\]
\end{frame}


\begin{frame}{Q4(c): Consistency and Order of Accuracy}

From the truncation error estimate,
\[
\mathcal{T}_{i,j}=O(h^2)+O(k).
\]
Thus,
\[
\mathcal{T}_{i,j}\longrightarrow 0
\qquad\text{as }h,k\longrightarrow0,
\]
so the scheme is consistent.
\begin{block}{Conclusion}
The scheme is consistent with order $(2,1)$: second-order accurate
in space and first-order accurate in time.
\end{block}

\end{frame}


\begin{frame}{Q4(d): Determine the Stability Condition}

To analyze stability, consider one Fourier mode
\[
U_i^j=\widehat U^{\,j}(\xi)e^{\mathrm{i}i\xi},
\qquad \mathrm{i}=\sqrt{-1}.
\]
A shift by one grid point gives
\[
U_{i+1}^j=e^{\mathrm{i}\xi}U_i^j,
\qquad
U_{i-1}^j=e^{-\mathrm{i}\xi}U_i^j.
\]

For the centered second difference,
\[
(BU^j)_i
=
\frac{U_{i-1}^j-2U_i^j+U_{i+1}^j}{h^2}
=
\underbrace{\frac{e^{-\mathrm{i}\xi}-2+e^{\mathrm{i}\xi}}{h^2}}_{\widehat B(\xi)}
U_i^j, \quad \widehat B(\xi)=-\frac{4}{h^2}\sin^2\frac{\xi}{2}.
\]
For the centered first difference,
\[
(CU^j)_i
=
\frac{U_{i+1}^j-U_{i-1}^j}{2h}
=
\underbrace{\frac{e^{\mathrm{i}\xi}-e^{-\mathrm{i}\xi}}{2h}}_{\widehat C(\xi)}
U_i^j,
\qquad
\widehat C(\xi)=\frac{\mathrm{i}\sin\xi}{h}.
\]

\end{frame}


\begin{frame}{Q4(d): Determine the Stability Condition (cont'd)}

The matrix scheme is
\[
\left(I-\frac{k}{2}B\right)U^{j+1}
=
\left(I+\frac{k}{2}B-kC\right)U^j.
\]

For a Fourier mode, $U_i^{j+1}=g(\xi)U_i^j$. Using the symbols of
$B$ and $C$,
\[
\left(1+2\frac{k}{h^2}\sin^2\frac{\xi}{2}\right)g(\xi)
=
1-2\frac{k}{h^2}\sin^2\frac{\xi}{2}
-\mathrm{i}\frac{k}{h}\sin\xi.
\]

Therefore, the amplification factor is
\[
g(\xi)=
\frac{1-2\frac{k}{h^2}\sin^2\frac{\xi}{2}
-\mathrm{i}\frac{k}{h}\sin\xi}
{1+2\frac{k}{h^2}\sin^2\frac{\xi}{2}}.
\]

\end{frame}


\begin{frame}{Q4(d): Determine the Stability Condition (cont'd)}

The scheme is stable if $|g(\xi)|\le1$ for all $\xi\in[-\pi,\pi]$.
Using the amplification factor,
\[
|g(\xi)|^2
=
\frac{
\left(1-2\frac{k}{h^2}\sin^2\frac{\xi}{2}\right)^2
+\left(\frac{k}{h}\sin\xi\right)^2
}{
\left(1+2\frac{k}{h^2}\sin^2\frac{\xi}{2}\right)^2
}.
\]

Thus, $|g(\xi)|\le1$ is equivalent to
\[
\left(\frac{k}{h}\sin\xi\right)^2
\le 8\frac{k}{h^2}\sin^2\frac{\xi}{2},
\]
which reduces to
\[
k\cos^2\frac{\xi}{2}\le2.
\]

Since $\max_{\xi\in[-\pi,\pi]}\cos^2(\xi/2)=1$, the stability condition is
\[
\boxed{k\le2.}
\]

\end{frame}


\section[Numerical]{Numerical Example}

\begin{frame}{Question 5}

Consider a European call option with strike $K$ and expiry $T$. The
underlying asset follows the geometric Brownian motion
\[
\dd S_t=\mu S_t\,\dd t+\sigma S_t\,\dd W_t,
\]
and the risk-free interest rate is $r$. Use the parameters
\[
S_0=K=100,\qquad r=0.05,\qquad \sigma=0.15,\qquad T=1.
\]

After transforming to log-price and time-to-expiry coordinates, the
Black--Scholes equation is
\[
V_\tau=\frac12\sigma^2V_{xx}+\nu V_x-rV,
\qquad
\nu=r-\frac12\sigma^2,
\qquad (x,\tau)\in\mathbb{R}\times[0,T].
\]
The initial condition is
\[
V(x,0)=(e^x-K)^+.
\]
\end{frame}

\begin{frame}{Q5(a): Local Truncation Error and Convergent Rate}

The scheme uses a forward difference in time and centered differences
in space:
\[
\frac{V_i^{j+1}-V_i^j}{k}
=
\frac{\sigma^2}{2}
\frac{V_{i-1}^j-2V_i^j+V_{i+1}^j}{h^2}
+\nu\frac{V_{i+1}^j-V_{i-1}^j}{2h}
-rV_i^{j+1}.
\]
Define
\[
\rho:=\frac{\sigma^2 k}{2h^2},\qquad
p_u:=\rho\left(1+\frac{\nu h}{\sigma^2}\right),\quad
p_0:=1-2\rho,\quad
p_d:=\rho\left(1-\frac{\nu h}{\sigma^2}\right).
\]
Then the scheme becomes
\[
V_i^{j+1}
=
\frac{p_dV_{i-1}^j+p_0V_i^j+p_uV_{i+1}^j}{1+rk}.
\]
    
\end{frame}

\begin{frame}{Q5(a): Local Truncation Error and Convergent Rate (cont'd)}

Substitute the exact solution into the scheme and write
\[
V_{i,j}:=V(x_i,\tau_j).
\]
Define the centered difference approximations
\[
\delta_x V_{i,j}
:=\frac{V_{i+1,j}-V_{i-1,j}}{2h},
\qquad
\delta_{xx}V_{i,j}
:=\frac{V_{i-1,j}-2V_{i,j}+V_{i+1,j}}{h^2}.
\]

The local truncation error at $(x_i,\tau_j)$ is
\[
\mathcal{T}_{i,j}
=
\frac{V_{i,j+1}-V_{i,j}}{k}
-\frac{\sigma^2}{2}\delta_{xx}V_{i,j}
-\nu\delta_xV_{i,j}
+rV_{i,j+1}.
\]

\end{frame}

\begin{frame}{Q5(a): Local Truncation Error and Convergent Rate (cont'd)}

Assume $V$ is sufficiently smooth for $\tau>0$, with the required
derivatives bounded. At $(x_i,\tau_j)$,
\[
\frac{V_{i,j+1}-V_{i,j}}{k}
=
V_\tau+\frac{k}{2}V_{\tau\tau}+O(k^2).
\]

The centered differences satisfy
\[
\delta_{xx}V_{i,j}=V_{xx}+O(h^2),
\qquad
\delta_xV_{i,j}=V_x+O(h^2).
\]

The discount term at the new time level satisfies
\[
V_{i,j+1}=V+kV_\tau+O(k^2).
\]

All derivatives are evaluated at $(x_i,\tau_j)$.

\end{frame}


\begin{frame}{Q5(a): Local Truncation Error and Convergent Rate (cont'd)}

All derivatives are evaluated at $(x_i,\tau_j)$. Substituting the
Taylor expansions into $\mathcal{T}_{i,j}$ gives
\[
\begin{aligned}
\mathcal{T}_{i,j}
={}&
\underbrace{V_\tau-\frac{\sigma^2}{2}V_{xx}-\nu V_x+rV}_{=\,0\text{ by the PDE}}\\
&+k\left(\frac12V_{\tau\tau}+rV_\tau\right)
+O(h^2)+O(k^2).
\end{aligned}
\]

Since $V$ satisfies the Black--Scholes PDE, the leading terms cancel.
Assuming the required derivatives are bounded,
\[
\boxed{\mathcal{T}_{i,j}=O(h^2)+O(k).}
\]

\end{frame}


\begin{frame}{Q5(b): Boundary Conditions and Grid Parameters}

For the truncated domain $[-1,5]\times[0,1]$, impose
\[
V(-1,\tau)=0,
\qquad
V(5,\tau)=e^5-Ke^{-r\tau}.
\]

With $h=0.01$ and $k=1/225$,
\[
\sigma\sqrt{k}=h,
\qquad
h<\frac{\sigma^2}{\nu}\approx0.58065.
\]
Thus, condition (3.9) is satisfied. The scheme coefficients are
\[
\rho=\frac12,\qquad
p_u\approx0.5086111,\quad
p_0=0,\quad
p_d\approx0.4913889.
\]

Starting from $V(x,0)=(e^x-K)^+$, advance the scheme to $\tau=1$
and interpolate at $x=\ln100$:
\[
\boxed{V(\ln100,1)\approx8.59764062.}
\]

\end{frame}


\begin{frame}{Q5(c): Enlarged Computational Domain}

Enlarge the domain to $[-2,10]\times[0,1]$. The boundary conditions are
\[
V(-2,\tau)=0,
\qquad
V(10,\tau)=e^{10}-Ke^{-r\tau}.
\]

Using the same $h=0.01$ and $k=1/225$, advance the scheme from
$V(x,0)=(e^x-K)^+$ to $\tau=1$, then interpolate at $x=\ln100$:
\[
\boxed{V(\ln100,1)\approx8.59763953.}
\]

The two domain choices give
\[
\left|8.59764062-8.59763953\right|
\approx1.10\times10^{-6}.
\]
Thus, for this grid, enlarging the computational domain has a negligible
effect on the computed option value.

\end{frame}




\begin{frame}\vfill\begin{center}{\LARGE Thank You! Questions?}\end{center}\vfill\end{frame}
\end{document}
